\pg{334} 56. O ente,. qué é acto puro, é necessario ; e o ente, que é necessario, é acto puro.

a) \emph{O ente}, \emph{que} é acto \emph{puro}, \emph{é necessario}. --- O .ente, que é: iE \emph{acto puro}, exclue-de-si-toda-a potencia : aliás não seria o que é. Ora o ente, que exclue .de si toda a potencia, deve possuir a existencia identificada .com a essencia, e por isso é \emph{necessario}. Logo o ente, que é \emph{acto} puro, é necessario \footnote{0 \emph{acto puro} exclue de si toda a potencia. Ora \emph{o} ênte, que exclue de si toda a potencia, não pode ser "contingente; porque o ente contingente não exclue, mas encerra -a potencia, emquanto está em potencia para o acto da existencia. Portanto'o ente, que é \emph{acto puro}, € \emph{necessario}, possue a existencia pela exigencia: da sua essencia. Nem podia ser: diversamente, O ente, que é \emph{acto} puro, possue a existencia identificada com a essencia, de modo que podemos dizer que a existencia é um dos elementos da essencia, Sendo assim, esse ente. é necessariamente \emph{existente}, como o homem, que pela sua essência é dotado: de razão,; É necessariamente racional. As essências, como \textasciicircum{} vimos, são necessárias quanto 208. elementos, que as constituem. .}.

b) \emph{O} ente, \emph{que} é necessario, é \emph{acto}. \emph{puro}. --- O ente, que é \emph{necessario}, existe pela exigencia da Sua essencia, e por isso. possue a existencia identificada com a essencia. Ora o ente, que possue a existencia identificada com. a essencia, é \emph{acto puro}. “Logo o: ente, que é

é necessario, é \emph{acto puro} \footnote{O ente \emph{necessario} é acto puro: O \emph{acto} é puro, quando não é recebido num sujeito distincto, nem está em potencia \textbf{para} um acto ulterior. Ora o ente necessario --- não é recebido “num sujeito distincto, porque o seu ser é uma e a mesma coisa com. a essencia, --- nem está em potencia para um acto ulterior, porque. o ente, ,eujo ser se identifica com a essencia, é o proprio ser, é o ser subsistente, e por isso possue a plenitude .da perfeição. .\par D' onde se vê que o ente necessario é unico. Por quanto o acto, como dissemos, só se multiplica emquanto é recebido em muitos sujeitos, Ora o ser do ente necessario, identificando-se com a essencia,.não é recebido em nenhum sujeito, mas é por si subsistente, e por isso não se multiplica, mas é unico. Logo o ente \emph{necessario} é unico.}. . ! 57. O ente, que consta de acto e de potencia, é contingente: eo ente, que é contingente, consta de acto e de potencia.

; .8) \emph{O ente}, que consta. \emph{de}; acto e \emph{de potencia}, é - \emph{contingente}. ---. \emph{AL} O ente, que consta" de \emph{acto} e. de \emph{potência}, possue a existencia realmente distiüetà da essencia. Ora o ente, que possue a exi- 'stencia \emph{réalinente}, distineta. da essencia, é \emph{contingente}, Logo o” ente, que consta de \emph{acto} e de “\emph{potencia}, \emph{é contingente e})\footnote{A conclusão é certissima, Na verdade, todo o ente composto pode ser dissolvido; e, podendo ser dissolvido, é naturalmente contingente. --- Além.d'isso, o \emph{ente} composto é \emph{posterior} ás partes componentes, e por isso não pode \emph{ser necessario}, mas deve ser \emph{contingente}: \emph{porque} o ente necessario,\pgn{335}sendo o proprio ser, não está ordenado para. outro. ser mais perfeito, e é simplesmente o \emph{primeiro}. --- Mais: -o composto. depende, no seu ser, de uma causa, que une as. parte componentes e que deve ser não só anterior ao proprio composto, mas distincta delle. Dependendo, no seu ser; de uma causa. distincta, o composto não existe pela sua natureza, e por isso é contingente.} 